\subsection{模型优点}

\begin{enumerate}
  \item 四问共用同一径向守恒框架，由常物性、变物性耦合到收缩变域逐步加机制，假设和方程可以逐条对应到题给附录，而不是更换一套互不衔接的算法名称。
  \item 空间用节点有限体积，边界通量有直接物理含义，便于检查热量与水分收支；时间用隐式 BDF，适合表面薄层引起的刚性。
  \item 第一问温度与 Bessel--Duhamel 解一致，冻结扩散系数问题可与级数解对照，从而把离散误差与物理简化分开讨论。
  \item 输出半径落在网格节点上，题目要求的四位小数表经过网格与时间加密后保持稳定。
\end{enumerate}

\subsection{模型缺点}

\begin{enumerate}
  \item 主模型不显式计入蒸发潜热。表面失水已经进入质量平衡，但汽化耗热未进入能量方程。若把给定密度解释为初始湿密度，则预热阶段隐含潜热与显热增量为同一量级，不能以``预热时间短''断言潜热可忽略\cite{dasilva2014}。题面缺少吸附等温线和相容的界面传质定义，目前无法在不引入额外参数的前提下把潜热闭合进主计算。
  \item 空气含湿量与药材干基含水率被当作同一传质势作差。同类圆柱干燥研究中的边界往往使用物料平衡含水率，而不是空气含湿量本身\cite{dasilva2014,kaya2007}。
  \item 一维径向近似忽略端面；长时干燥中该误差会积累。
  \item 第四问的均匀收缩与有效热密度是在仅有外半径轨迹时的闭合办法，不是已经由内部变形观测证实的力学模型。
\end{enumerate}

\subsection{推广}

同一套有限体积--BDF 框架可以推广到：在能量边界中加入与失水通量一致的潜热项；用轴对称二维模型估计端部影响；在获得吸附等温线后把 $C_a$ 换成平衡含水率；以及把已知 $R(t)$ 换成由局部含水率驱动的收缩律。推广时必须补充相应参数，并重新做守恒与加密检验，不能直接套用其他农产品的经验系数。
